This thesis has examined the compact tail-assignment formulation from the perspectives of mathematical correctness, operational fidelity, and computational performance. The work proceeds from a systematic identification of two structural deficiencies in the original formulation of \citet{khaled2018compact} to a corrected and extended model that is both operationally valid and computationally tractable. All empirical claims in this thesis are limited to the project’s synthetic feasible-instance families and to the declared solver protocol; they are not presented as universal statements about all real airline schedules or all operational environments.

\paragraph{Correction of the basic feasibility model.}
The standard balance constraints of the compact formulation admit simultaneous-departure and temporal-overlap configurations that are physically impossible for a single aircraft. The corrected model appends pairwise conflict-exclusion constraints indexed by the set $\mathcal{O}$ of overlapping flight pairs, as defined in Chapter~\ref{ch:ground}, thereby restricting the feasible region to trajectories that are executable under turn-time requirements. This correction is a prerequisite for any maintenance model built on top of the assignment layer: without it, the maintenance constraints are evaluated over a distorted domain.

\paragraph{Correction of Constraint~(13).}
The central contribution of this thesis is the identification and correction of the logical flaw in Constraint~(13), the cumulative flight-hour bound. The original single-constraint formulation is shown to be vacuously non-binding when one endpoint maintenance indicator is inactive, permitting unchecked flight-hour accumulation in violation of the regulatory intent. The corrected formulation replaces the single constraint with the endpoint-anchored pair~\eqref{eq:c13a}--\eqref{eq:c13b}, each binding at the corresponding interval endpoint. The corrected pair is provably tighter: the intersection of the two anchored half-spaces is at most as permissive as the original, and strictly tighter on any maintenance-active instance where at least one endpoint indicator is zero at the LP optimum. The computational study of Chapter~\ref{ch:experiments} validates the correction directly: on the dedicated loophole instance the original form returns an optimal solution that an independent validator shows to violate the flight-hour limit by 400 minutes, whereas the corrected form is proved infeasible. The associated LP-relaxation tightening is quantified separately in the strengthening study of Chapter~\ref{ch:strengthening}.

\paragraph{Full A/B/C/D hierarchy extension.}
The maintenance model is extended to the four-level regulatory hierarchy in Chapter~\ref{ch:hierarchy}, incorporating type-specific flight-hour thresholds, calendar-day spacing constraints, multi-day maintenance windows, the subsumption ordering $D \succ C \succ B \succ A$, and pre-horizon state accounting. The extension preserves the compact assignment structure while capturing the full regulatory framework governing airline maintenance operations.

\paragraph{Heuristic methods and computational evaluation.}
Chapter~\ref{ch:heuristics} presents seven executable heuristic modes --- greedy, insertion, greedy-plus-insertion, repair, bounded local search, a modified Dijkstra network search, and ant-colony optimisation --- all of which invoke a common timeline-feasibility oracle so that a valid route has the same meaning across methods. The computational study of Chapter~\ref{ch:experiments} demonstrates that exact optimisation is tractable on the small and medium exact-feasible instances, and that the bounded local search improves the greedy objective by a mean of 12.7\% on the 25-instance matrix. Consistent with the coverage-versus-cost distinction developed in that chapter, heuristic-to-MILP differences are reported descriptively rather than as universal optimality gaps, because the heuristic timeline and the integrated MILP do not enumerate identical maintenance and ferry decisions on every instance.

\paragraph{MILP strengthening and current CPLEX validation.}
Chapter~\ref{ch:strengthening} studies opt-in formulation-strengthening methods --- aircraft-compatible maintenance domains, deferred-trigger reachability, strong trigger linking, interval-specific big-$M$ values, and pairwise maintenance conflicts --- while keeping the corrected baseline unchanged by default. On a family of 20 constructive, timeline-validated feasible instances solved with the commercial CPLEX 12.6 AMPL driver, the combined formulation preserves every certified MILP objective, raises the LP lower bound on all matched instances, and reduces mean exact-solve time by 15.8\%. This is the strongest computational evidence in the thesis for the strengthening direction, and it is explicitly presented as validated family evidence under the tested A-check family and solver protocol rather than evidence of universal superiority across all instance classes or real airline schedules.

\section*{Directions for future research}

Several directions merit further investigation. First, the corrected compact model could be strengthened with additional valid inequalities derived from the structure of the maintenance hierarchy, potentially yielding tighter LP bounds and faster exact convergence on high-pressure instances. Second, integration with fleet-assignment or disruption-recovery decisions would extend the model to a more complete operational planning context. Third, decomposition methods such as Benders decomposition or Lagrangian relaxation could be applied to scale the exact formulation to larger instance families without sacrificing optimality guarantees. Finally, the heuristic framework could be augmented with adaptive or population-based meta-heuristics and evaluated on instances derived from real airline schedules to assess practical transferability beyond the parameterised benchmark suite.
